\section{Discussion}
\label{sec:discussion}

The theoretical part of the paper makes one distinction precise. With fixed, known reporting, exogenous exposure and conditionally independent counting errors, the risk for the current value is $A/m$ and the risk for the realised mean of the observed periods is $A/(Tm)$. The process mean adds the variance of the realised mean, which more cases within the same periods do not reduce. The count at which counting error equals between-period variance is the half-weight point of the corresponding shrinkage estimator, as in small-area estimation and disease mapping \citep{fay1979estimates, clayton1987empirical}. It belongs with its estimator, target and sampling experiment and is not a property of a pathogen. The efficiency statement holds within a class of conditionally unbiased estimators with a known denominator. Information from other features of the epidemic curve \citep{parag2022quantifying, parag2022fundamental}, independent reporting measurements or prior structure changes what can be learned, and estimators that accept bias in exchange for lower variance lie outside that class. Appendix~\ref{app:extensions} gives the corresponding forecast weight, the allocation of a fixed budget between periods and cases, and a lower bound on the precision with which a few seasons can estimate between-season variation.

The empirical part shows what can and cannot be read from a cross-sectional dispersion curve. Its intercept is a count-independent relative dispersion, the part of the variation in the statistic that does not scale with the count. With reporting fixed across weeks it is the relative variance of transmission across units and weeks; with week-to-week reporting shocks it also contains their contribution, and single-period moments cannot split the two (Proposition~\ref{prop:report-noise}). For influenza admissions it is about 0.015. Under pure transmission heterogeneity, scanning the entire $(\rho, k)$ parameter manifold compatible with the counting slope gives a tightly bound transmission relative standard deviation between 11.6\% and 12.3\% (9.9\% and 13.8\% propagated); expanding across the full grid of reporting shocks and scale dependence yields a global envelope between 5.0\% and 20.4\%. Under fixed reporting this is the variation that a current-value analysis wants to resolve and that a process-mean analysis treats as sampling variation; under volatile reporting it is an upper bound on the error floor of the current value. The counting slope, its ratio to the independent-Poisson or branching-Poisson reference and the thinning response cannot distinguish overdispersed offspring, under-reporting and heterogeneous transmission. The audit produces slopes near the observed one under several combinations, and the grid over $\rho$ and $k$ admits values from heavy under-reporting with strong superspreading to near-complete reporting with moderate dispersion. The dispersion-to-counting crossover count was about 37 admissions per week for the pooled panel. Equal season intercepts were rejected under both simulated null distributions and non-parametric data calibrations (Wild cluster bootstrap $p=0.013$; permutation $p=0.045$), mainly because of a lower intercept in 2023--24, confirming seasonal transmission divergence across surveillance cycles. The crossover marks where counting noise stops dominating the cross-sectional dispersion of this statistic, not where additional cases stop being useful for a particular target.

\subsection{Real-time estimation}
\label{sec:disc-realtime}

All analyses used a retrospective snapshot in which counts are close to final. In real time, reporting delays leave the most recent weeks incomplete, and right-censoring biases the latest counts downwards in proportion to the delay distribution \citep{hohle2014bayesian, mcgough2020nowcasting}. The current-value target is affected first and most, because its risk is dominated by the last one or two periods, exactly those that are incomplete, whereas a realised or process mean over many periods dilutes them. The $R_A$ statistic also needs $C_{t+1}$, so it is available only one week after its reference week, and estimates near the end of a series combine smaller effective counts with nowcasting uncertainty. Real-time renewal methods address this by modelling delays and the infection process jointly \citep{abbott2020estimating, lison2024generative, gostic2020practical}. In the present framework, nowcasting adds a correction-error term to the counting term, and its variance grows as the delay distribution places more mass beyond the current week. A crossover count estimated from retrospective data will therefore overstate the precision available in real time.

Estimation windows raise a related issue. Widening a window in renewal estimators accumulates incidence while imposing constancy of the reproduction number over a longer interval \citep{cori2013new, gostic2020practical}. This trades temporal resolution for precision and is not equivalent to observing more independent periods; the averages in Table~\ref{tab:targets} instead give each period its own transmission value. Applying their formulas to overlapping windows requires explicit treatment of the overlap.

\subsection{Reporting that depends on transmission}
\label{sec:disc-reporting}

Proposition~\ref{prop:report-noise} assumes that the reporting probability is independent of transmission. In practice testing, care seeking and admission thresholds can change with epidemic activity \citep{pitzer2021impact, sherratt2021exploring}. Write $z_t=p_t/\rho$ for the relative reporting probability. Given $R_t$, the estimator then has conditional bias
\begin{equation}
  \E(\Rhat_t\mid R_t)-R_t=\left[\E(z_t\mid R_t)-1\right]R_t .
  \label{eq:correlated-bias}
\end{equation}
By Eq.~\eqref{eq:correlated-bias}, if reporting rises with transmission, so that $\Cov(z_t,R_t)>0$, high-transmission periods are over-counted and low-transmission periods under-counted, the estimated swings exaggerate the true ones, and the fitted intercept overstates the relative dispersion of transmission. If reporting falls with transmission, for example when admission capacity saturates or testing is rationed at the peak, the estimates damp the swings and the intercept understates it. In both cases the process-mean estimate is biased by $\Cov(z_t,R_t)$, and the bias does not shrink with more cases. Neither direction can be detected from the counts alone; it requires external information on testing or admission practice. The participation adjustment used here addresses only facility participation, and it changed the intercept negligibly.

\subsection{Choice of interval distribution}
\label{sec:disc-interval}

The transform in Eq.~\eqref{eq:ra} needs an interval distribution. We kept the household serial interval of \citet{cowling2009estimation} as the primary choice because it is directly estimated for influenza, widely used and fixed before the analyses reported here. The intrinsic generation-time estimate of \citet{chan2025estimating}, the more appropriate quantity for a renewal mapping, served as a sensitivity analysis. Serial intervals, intrinsic generation intervals and realised household generation intervals are different quantities \citep{vink2014serial}, and the mapping from growth rate to reproduction number depends on which is used \citep{wallinga2007how}. The levels of $a$ and $b$ change with the interval, but their ratio and the slope ratio to the independent-Poisson reference change by less than their sampling uncertainty, so the conclusions about what the diagnostics measure do not depend on this choice.

\subsection{Identification and limitations}
\label{sec:disc-limits}

Separating the mechanisms behind the dispersion curve requires information that a single cross-section lacks. Temporal structure can separate persistent transmission variation from independent counting and reporting noise under explicit state-space or renewal models \citep{koyama2021estimating, knape2008estimability, augermethe2016state}, and external measurements of reporting can constrain $\rho$ and its variation. The shape of the curve narrows the candidates without identifying one. The observed power-minus-affine contrast lay at or below the 0.3\% quantile of every generator whose heterogeneity does not depend on count (S1 and S3--S6). Scale-dependent heterogeneity alone (S7a) reproduces the contrast but not the slope ratio. Scale-dependent heterogeneity combined with reporting shocks (S8) reproduces the intercept, the crossover count, the slope ratio and the contrast, and the share grid shows a family of such combinations. Within that family the intercept is between 1.04 and 3.85 times the transmission relative variance, so the share of $a$ due to transmission remains unidentified. The held-out state contrast gives the same ordering: it is improbable under count-independent heterogeneity and lies in the lower tail of S8 and S9. No mechanism in the audit or the grid reproduces the gap between pooled and within-state slopes, and week centring is reproduced only where part of the heterogeneity is national, so the family is an approximation and not a description of the data-generating process. Reporting that varies with jurisdiction size would leave the same trace.

The growth windows are a fixed rule applied to four seasons of one surveillance system, and inference conditions on that selection; window thresholds of 1.02 and 1.10 left the crossover count at 37.4 and 38.6, and transfer of these values to other epidemic regimes would need data from those regimes. The analyses used a retrospective snapshot, so reporting delays and right-censoring of recent weeks were not modelled (Section~\ref{sec:disc-realtime}). The audit generators are stylised, and their count alignment matched the median and lower tail of the observed counts more closely than the upper tail, although reweighting to the observed count deciles changed the intercept recovery by at most 3.2\%. During major epidemic surges where counts enter the upper tail ($C > 389$), potential non-linear admission thresholds or capacity constraints could induce additional scale-dependent compression of apparent dispersion, underscoring the value of auditing across flexible tail geometries. The pigeonhole bootstrap intervals were conservative in the audit, with standard deviations 1.6--2.5 times the Monte Carlo values. The two-way sandwich intervals reported for $a$, $b$ and $b/a$ met the coverage rule, but only at its threshold under S6 and S8, and intervals for other quantities remain bootstrap intervals. While the parametric generators have less dependent variation than the data, non-parametric Wild cluster bootstrap and within-state permutation tests directly on the empirical data independently corroborate the seasonal divergence at the 5\% level. Scanning the full $(\rho, k)$ manifold confirms that the implied transmission relative standard deviation remains bounded within 9.9--13.8\% across all slope-compatible cells under pure heterogeneity, and within 5.0--20.4\% across mechanisms with reporting shocks. Intervals for iterated fits condition on convergence, and within-unit slopes were imprecise and sensitive to the weighting scheme. Facility participation measures only one aspect of ascertainment, and the comparison with a second reporting regime in Appendix~\ref{app:china} is descriptive.

\subsection{Checklist for surveillance practice}
\label{sec:disc-checklist}

Four recommendations for reporting reproduction numbers follow from these results.
\begin{enumerate}[leftmargin=*,itemsep=2pt]
\item State the inferential target. An estimate may target the current value $R_t$, the realised mean $\Rbar_T$ over the observed periods or the process mean $\Rbar$. Precision statements of the form $A/m$ apply to the first two under declared reporting assumptions and omit between-period variation for the third.
\item Inflate the counting error when the denominator is an observed count. The delta factor is $1+(1-\rho)(\Rbar^2+v_R)/A$, which is 1.48 under the illustrative parameters of Section~\ref{sec:model-observed}. At small counts the factor is larger, 2.42 at $m=10$ in Table~\ref{tab:ma_targets_observed}. The half-weight count of the shrinkage estimator moves from $\mx$ to $\mx+(1-\rho)(\Rbar^2+v_R)/v_R$, from 53.1 to 78.6 under the same parameters.
\item Qualify single-period estimates for reporting volatility. One period's counts cannot separate transmission changes from reporting shocks (Proposition~\ref{prop:report-noise}), so rapid changes should be checked against testing volumes, reporting participation or temporal models before they are interpreted as transmission.
\item Treat dispersion fits as descriptive. The intercept $a$, the slope $b$ and the crossover count $b/a$ describe the observed statistic on the selected rows. They are not a transmission-variance estimate without an assumed mechanism, and $b/a$ is not a sample-size threshold for surveillance.
\end{enumerate}
